\section{Вычитание целых чисел}
\subsection{Вычитание чисел A и В c различными знаками}

В таблице~\ref{tab:a-b-dop2} представлены прямые и дополнительные коды
для чисел $A$ и $B$.

\begin{table}[H]
    \centering
    \caption{Прямые и дополнительные коды для A и В}
	\label{tab:a-b-dop2}
    \begin{tabular}{|c|c|c|}
		\hline
		 & Прямой код & Дополнительный код \\
		\hline
		$A=109$ & $0110\ 1101$ & \makecell{$2^8-109$\\$1001\ 0011$}  \\
		\hline
		$B=17$ & $0001\ 0001$ & \makecell{$2^8-17$\\$1110\ 1111$} \\
		\hline
    \end{tabular}
\end{table}

\subsubsection{A-B}

\begin{table}[H]
    \centering
    \caption{Вычитание А и В}
    \begin{tabular}{c|cc|cccccccp{1cm}ccp{0.5cm}cc}
		\small заём & & & & $\bullet$ & $2$ & & & & & & \multicolumn{2}{c}{знаковая} & \  & \multicolumn{2}{c}{беззнаковая}  \\
		\cline{1-10} \cline{12-13} \cline{15-16}
		$A_{\text{пр}}$ & & 0 & 1 & 1 & 0 & 1 & 1 & 0 & 1 & & \multirow{2}{*}{-} & 109 & \  & \multirow{2}{*}{-} & 109   \\
		$B_{\text{пр}}$ & & 0 & 0 & 0 & 1 & 0 & 0 & 0 & 1 & & & 17 & \  &  & 17 \\
		\cline{1-10} \cline{12-13} \cline{15-16}
		$C_{\text{пр}}$ & & 0 & 1 & 0 & 1 & 1 & 1 & 0 & 0 & & & 92 & \  &  & 92 \\
    \end{tabular}
\end{table}

\begin{table}[H]
    \centering
    \begin{tabular}{cccccc}
        CF=0; & ZF=0; & PF=1; & AF=0; & SF=0; & OF=0;
    \end{tabular}
\end{table}

\subsubsection{-A-B}

\begin{table}[H]
    \centering
    \caption{Вычитание -А и В}
    \begin{tabular}{c|cc|cccccccp{1cm}ccp{0.5cm}cc}
		\small заём & & & & & & & & & & & \multicolumn{2}{c}{знаковая} & \  & \multicolumn{2}{c}{беззнаковая}  \\
		\cline{1-10} \cline{12-13} \cline{15-16}
		$A_{\text{доп}}$ & & 1 & 0 & 0 & 1 & 0 & 0 & 1 & 1 & & \multirow{2}{*}{-} & -109 & \  & \multirow{2}{*}{-} & 147  \\
		$B_{\text{пр}}$ & & 0 & 0 & 0 & 1 & 0 & 0 & 0 & 1 & & & 17 & \  &  & 17 \\
		\cline{1-10} \cline{12-13} \cline{15-16}
		$C_{\text{доп}}$ & & 1 & 0 & 0 & 0 & 0 & 0 & 1 & 0 & & & & \  &  & 130 \\
		\cline{1-10} \cline{12-13} \cline{15-16}
		$C_{\text{пр}}$ & & 1 & 1 & 1 & 1 & 1 & 1 & 1 & 0 & & & -126 & \  &  & 
    \end{tabular}
\end{table}

\begin{table}[H]
    \centering
    \begin{tabular}{cccccc}
        CF=0; & ZF=0; & PF=0; & AF=0; & SF=1; & OF=0;
    \end{tabular}
\end{table}


\subsubsection{A-(-B)}

\begin{table}[H]
    \centering
    \caption{Вычитание А и -В}
    \begin{tabular}{c|cc|cccccccp{1cm}ccp{0.5cm}cc}
		\small заём & $\bullet$ & $1$ & $1$ & $1$ & $1$ & $1$ & $1$ & $2$ & & & \multicolumn{2}{c}{знаковая} & \  & \multicolumn{2}{c}{беззнаковая}  \\
		\cline{1-10} \cline{12-13} \cline{15-16}
		$A_{\text{пр}}$ & & 0 & 1 & 1 & 0 & 1 & 1 & 0 & 1 & & \multirow{2}{*}{-} & 109 & \  & \multirow{2}{*}{-} & 109   \\
		$B_{\text{доп}}$ & & 1 & 1 & 1 & 0 & 1 & 1 & 1 & 1 & & & -17 & \  &  & 239 \\
		\cline{1-10} \cline{12-13} \cline{15-16}
		$C_{\text{пр}}$ & & 0 & 1 & 1 & 1 & 1 & 1 & 1 & 0 & & & 126 & \  &  & 126? \\
    \end{tabular}
\end{table}

\begin{table}[H]
    \centering
    \begin{tabular}{cccccc}
        CF=1; & ZF=0; & PF=1; & AF=1; & SF=0; & OF=0;
    \end{tabular}
\end{table}

Для беззнаковой интерпретации результат неверен из-за заёма разряда за пределами формата.

\subsubsection{-A-(-B)}

\begin{table}[H]
    \centering
    \caption{Вычитание -А и -В}
    \begin{tabular}{c|cc|cccccccp{1cm}ccp{0.5cm}cc}
		\small заём & $\bullet$ & $1$ & $1$ & $2$ & $\bullet$ & $1$ & $2$ & & & & \multicolumn{2}{c}{знаковая} & \  & \multicolumn{2}{c}{беззнаковая}  \\
		\cline{1-10} \cline{12-13} \cline{15-16}
		$A_{\text{доп}}$ & & 1 & 0 & 0 & 1 & 0 & 0 & 1 & 1 & & \multirow{2}{*}{-} & -109 & \  & \multirow{2}{*}{-} & 147\\
		$B_{\text{доп}}$ & & 1 & 1 & 1 & 0 & 1 & 1 & 1 & 1 & & & -17 & \  &  & 239 \\
		\cline{1-10} \cline{12-13} \cline{15-16}
		$C_{\text{доп}}$ & & 1 & 0 & 1 & 0 & 0 & 1 & 0 & 0 & & & & \  &  & 164? \\
		\cline{1-10} \cline{12-13} \cline{15-16}
		$C_{\text{пр}}$ & & 1 & 1 & 0 & 1 & 1 & 1 & 0 & 0 & & & -92 & \  &  & \\
    \end{tabular}
\end{table}

\begin{table}[H]
    \centering
    \begin{tabular}{cccccc}
        CF=1; & ZF=0; & PF=0; & AF=1; & SF=1; & OF=0;
    \end{tabular}
\end{table}

Для беззнаковой интерпретации результат неверен из-за заёма разряда за пределами формата.

\subsection{Подбор числа В}
Для того, чтобы получить переполнение, необходимо, чтобы $A+B>128$. Возьмём
$B=20_{10}=0001\ 0100_{2}$. Тогда $-B=2^8-20=236_{10}=1110\ 1100_{2}$.

\begin{table}[H]
    \centering
    \caption{Вычитание -А и В}
    \begin{tabular}{c|cc|cccccccp{1cm}ccp{0.5cm}cc}
		\small заём & & $\bullet$ & $1$ & $1$ & $1$ & $1$ & $2$ & & & & \multicolumn{2}{c}{знаковая} & \  & \multicolumn{2}{c}{беззнаковая}  \\
		\cline{1-10} \cline{12-13} \cline{15-16}
		$A_{\text{доп}}$ & & 1 & 0 & 0 & 1 & 0 & 0 & 1 & 1 & & \multirow{2}{*}{-} & -109 & \  & \multirow{2}{*}{-} & 147   \\
		$B_{\text{пр}}$ & & 0 & 0 & 0 & 1 & 0 & 1 & 0 & 0 & & & 20 & \  &  & 20 \\
		\cline{1-10} \cline{12-13} \cline{15-16}
		$C_{\text{пр}}$ & & 0 & 1 & 1 & 1 & 1 & 1 & 1 & 1 & & & 127? & \  &  & 127 \\
    \end{tabular}
\end{table}

\begin{table}[H]
    \centering
    \begin{tabular}{cccccc}
        CF=0; & ZF=0; & PF=1; & AF=1; & SF=1; & OF=1;
    \end{tabular}
\end{table}

Для знакового вычитания результат неверный из-за переполнения формата.

\begin{table}[H]
    \centering
    \caption{Вычитание А и -В}
    \begin{tabular}{c|cc|cccccccp{1cm}ccp{0.5cm}cc}
		\small заём & $\bullet$ & $2$ &  &  &  &  & & & & & \multicolumn{2}{c}{знаковая} & \  & \multicolumn{2}{c}{беззнаковая}  \\
		\cline{1-10} \cline{12-13} \cline{15-16}
		$A_{\text{пр}}$ & & 0 & 1 & 1 & 0 & 1 & 1 & 0 & 1 & & \multirow{2}{*}{-} & 109 & \  & \multirow{2}{*}{-} & 109\\
		$B_{\text{доп}}$ & & 1 & 1 & 1 & 0 & 1 & 1 & 0 & 0 & & & -20 & \  &  & 236 \\
		\cline{1-10} \cline{12-13} \cline{15-16}
		$C_{\text{доп}}$ & & 1 & 0 & 0 & 0 & 0 & 0 & 0 & 1 & & & & \  &  & 129?\\
		\cline{1-10} \cline{12-13} \cline{15-16}
		$C_{\text{пр}}$ & & 1 & 1 & 1 & 1 & 1 & 1 & 1 & 1 & & & -127? & \  &  & \\
    \end{tabular}
\end{table}

\begin{table}[H]
    \centering
    \begin{tabular}{cccccc}
        CF=1; & ZF=0; & PF=0; & AF=0; & SF=0; & OF=1;
    \end{tabular}
\end{table}

Для знакового вычитания результат неверный из-за переполнения формата.
Для беззнакового вычитания результат неверный из-за заёма разряда за пределами формата. 

\subsection{Подбор числа А}
Для того чтобы при $A-(-B)$ давало переполнение, а при $-A-B$
результат был корректен, необходимо $A+B=128$. Значит, возьмём
$A=111_{10}=0110\ 1111_{2}$, тогда $-A=2^8-109=147_{10}=1001\ 0001_{2}$.

\begin{table}[H]
    \centering
    \caption{Вычитание А и -В}
    \begin{tabular}{c|cc|cccccccp{1cm}ccp{0.5cm}cc}
		\small заём & $\bullet$ & $2$ & & & & & & & & & \multicolumn{2}{c}{знаковая} & \  & \multicolumn{2}{c}{беззнаковая}  \\
		\cline{1-10} \cline{12-13} \cline{15-16}
		$A_{\text{пр}}$ & & 0 & 1 & 1 & 0 & 1 & 1 & 1 & 1 & & \multirow{2}{*}{-} & 111 & \  & \multirow{2}{*}{-} & 111   \\
		$B_{\text{доп}}$ & & 1 & 1 & 1 & 0 & 1 & 1 & 1 & 1 & & & -17 & \  &  & 239\\
		\cline{1-10} \cline{12-13} \cline{15-16}
		$C_{\text{доп}}$ & & 1 & 0 & 0 & 0 & 0 & 0 & 0 & 0 & & & & \  &  & -128? \\
		\cline{1-10} \cline{12-13} \cline{15-16}
		$C_{\text{пр}}$ & & 1 & 0 & 0 & 0 & 0 & 0 & 0 & 0 & & & -128? & \  &  & \\
    \end{tabular}
\end{table}

\begin{table}[H]
    \centering
    \begin{tabular}{cccccc}
        CF=1; & ZF=0; & PF=0; & AF=0; & SF=1; & OF=1;
    \end{tabular}
\end{table}

Для знакового вычитания результат неверен из-за переполнения формата.
Для беззнакового вычитания результат неверен из-за заёма разряда за пределами формата.

\begin{table}[H]
    \centering
    \caption{Вычитание -А и В}
    \begin{tabular}{c|cc|cccccccp{1cm}ccp{0.5cm}cc}
		\small заём & & & & & & & & & & & \multicolumn{2}{c}{знаковая} & \  & \multicolumn{2}{c}{беззнаковая}  \\
		\cline{1-10} \cline{12-13} \cline{15-16}
		$A_{\text{доп}}$ & & 1 & 0 & 0 & 1 & 0 & 0 & 0 & 1 & & \multirow{2}{*}{-} & -111 & \  & \multirow{2}{*}{-} & 145   \\
		$B_{\text{пр}}$ & & 0 & 0 & 0 & 1 & 0 & 0 & 0 & 1 & & & 17 & \  &  & 17\\
		\cline{1-10} \cline{12-13} \cline{15-16}
		$C_{\text{доп}}$ & & 1 & 0 & 0 & 0 & 0 & 0 & 0 & 0 & & & & \  &  & 128 \\
		\cline{1-10} \cline{12-13} \cline{15-16}
		$C_{\text{пр}}$ & & 1 & 0 & 0 & 0 & 0 & 0 & 0 & 0 & & & -128 & \  &  & \\
    \end{tabular}
\end{table}

\begin{table}[H]
    \centering
    \begin{tabular}{cccccc}
        CF=0; & ZF=0; & PF=0; & AF=0; & SF=1; & OF=0;
    \end{tabular}
\end{table}

